% =====================================================================
% arXiv v2 — BACK SECTIONS: §7 Deployment Evidence + §8 Discussion + §9 Conclusion
% REFRAME 2026-06-11 (Lars): CEP = optional advisory convention, not binding
% network enforcement. "activation/enforce" at the network scope -> "the
% convention's recomputable maturity signal"; CEP compels no one. Conformance
% stays "4/5 + 1 open item". Cross-protocol interop folded into §8.
% (a)/(b) substance + mechanics + citations + labels unchanged.
% =====================================================================

% ---------------------------------------------------------------------
\section{Deployment Evidence}
\label{sec:evidence}
% Material: v1.9 §4 + §6 re-weighted as SUPPORTING evidence (These-a).
% ---------------------------------------------------------------------

The governance design of Sections~\ref{sec:problem}–\ref{sec:scope} is worth posing only if the substrate it builds on is real. This section reports the deployment as supporting evidence — that the advisory verification layer runs, at scale, with independently checkable conformance — not as the paper's central claim.

\subsection{Live deployment, advisory layer running}

The reference implementation is live at \texttt{api.moltrust.ch} and has operated since March 2026. Eight credential verticals are operational (Core Identity, Commerce, Travel, Skill Verification, Prediction Markets, Brand Protection, Music, Sports Integrity), exposed through MCP tooling and a published SDK. Agent registrations, credential issuances, and confirmed violations are anchored on Base L2; an ERC-8004 ecosystem scanner indexes 44{,}355 agents for cross-ecosystem visibility independent of MolTrust registration \cite{ref37_erc8004}. The three layers of Section~\ref{sec:background} are live, and they run \textbf{advisory}: verdicts are signed, anchorable, and recorded; nothing is blocked at the network scope, and acting on a verdict is the relying party's decision. The deployment is itself a relying party of the protocol it defines, the natural first site of operator-local (scope (a)) enforcement.

\subsection{Scale and conformance}

As of the reporting window the registry holds a modest population — on the order of dozens of agents and endorsements (Section~\ref{subsec:sybil}) — against an external benchmark of 69{,}000 bots on a single marketplace \cite{ref1_agentmarket}: the substrate is real and the adversarial-scale validation is explicitly pending. Protocol conformance is documented in a published, SHA-256-anchored \texttt{CONFORMANCE.md} with Layer~A (protocol) and Layer~B (registry) requirements, auto-generated from source with drift detection, and a nine-check SKILL AUDIT mapped to CWE identifiers (cross-referencing NIST NVD and the Ferrag et al.\ taxonomy \cite{ref18_ferrag}). Independent verifiers can reproduce any check against the anchored specification using only a SHA-256 implementation and an HTTP client.

Against the Agent Interoperability Profile (AIP/IBCT) feature set \cite{ref27_kroehlssrn}, the implementation satisfies \textbf{four of the five} features, with the fifth — expressive chained policy — recorded as an \textbf{open item, not a pass}. The four implemented features hold against an independent verifier; the open item is expressive, multi-condition, context-dependent policy (Datalog/Biscuit/Cedar-class \cite{ref34_biscuit}). The current AAE constraint model uses URI-pattern matching, which is a deliberate lower bound on a single axis (Section~\ref{subsec:limits}): sufficient for the straightforward allow/deny boundaries of today's relying parties, but not equivalent to a policy language with nested logical and multi-party temporal conditions. We state this as ``4/5 + 1 open item'' rather than full conformance, and track the open item as roadmap rather than as a completed claim.

\subsection{Sybil resistance and behavioural consistency}
\label{subsec:sybil}

The substrate's defence against identity fabrication is the same machinery the CEP cluster-diversity condition reuses (Section~\ref{subsec:clusters}), which is why we report it here. Three mechanisms compound. \textbf{Dual-signature Interaction Proofs} require two distinct Ed25519 keys, turning Sybil forgery from cheap identity fabrication into expensive key compromise. \textbf{Cross-vertical diversity gating} applies a hard penalty when a non-seed agent draws endorsements from fewer than three distinct verticals, forcing a Sybil operator to maintain credible presence across unrelated domains. \textbf{Principal-DID-linked violation persistence} carries an operator's negative history across agent re-registration, defeating identity rotation. Cluster detection uses the reciprocal-Jaccard coefficient over endorser sets with a 0.8 production threshold — the identical edge predicate CEP's cluster count is built on (Section~\ref{subsec:clusters}), so the diversity the maturity condition demands is computed from the same structural signal the registry already trusts for Sybil defence. The economic layer (per-verify pricing) compounds rather than stands alone. These thresholds are operational heuristics calibrated at bootstrap scale; their behaviour under adversarial load is an open empirical question (Section~\ref{subsec:limits}).

% ---------------------------------------------------------------------
\section{Discussion}
\label{sec:discussion}
% Material: v1.9 §5 (cross-protocol interop, FOLDED here) + §7.2/§7.3/§7.4 +
% §17 open items. Positioning vs DAO/agent-ID/AIP; future work = §6.3.
% ---------------------------------------------------------------------

\subsection{Positioning against adjacent work}

CEP's object — an \emph{optional} network-wide trust convention whose maturity is recognised without a central authority — sets it apart from adjacent systems that either presuppose a central enforcer or coordinate nothing across organisations. Centralised architectures such as SAGA \cite{ref16_syros} obtain tight, formally analysed enforcement by trusting a Provider; CEP targets the case where no such Provider exists, offering a voluntary convention rather than per-domain enforcement. On-chain agent registries such as ERC-8004 \cite{ref37_erc8004} establish that an identity exists but not whether a network-wide trust convention around it has matured; CEP supplies that recomputable maturity signal. DAO governance frameworks \cite{ref39_daogov} decentralise \emph{treasury and parameter} decisions but typically concentrate voting power; CEP's concentration caps and five-condition conjunction (Section~\ref{subsec:params}) are stricter than the norms those analyses report, and gate a signal rather than a sanction. Against the AIP/IBCT profile \cite{ref27_kroehlssrn}, our position is the honest one of Section~\ref{sec:evidence}: four features met, expressive chained policy open.

\textbf{Cross-protocol interoperability} reinforces the positioning. A single Ed25519 + RFC~8785 signing operation produces artifacts verifiable by conformant implementations of independent protocols (verified against external Authority-Constraints and Provider-Attestation verifiers), with reproducible test vectors. Interoperability here is by shared cryptographic primitive rather than bilateral specification agreement, which is what lets the verdict of Section~\ref{sec:background} be consumed by parties that never integrated with MolTrust — a precondition for both operator-local (scope (a)) adoption and voluntary convergence on the convention without central coordination.

\subsection{Limitations}
\label{subsec:limits}

The limitations are explicit and bound the claims. \textbf{Scale:} the registry operates at bootstrap scale with a short production history; the Sybil heuristics (Jaccard 0.8, three-vertical minimum, penalty multiplier) are calibrated, not formally optimal, and untested under adversarial load — formal sensitivity analysis is planned at larger N. \textbf{Expressiveness:} URI-pattern matching is a deliberate lower bound; expressive chained policy is the AIP open item, not a delivered feature. \textbf{Kernel layer:} Falco eBPF detection is defence-in-depth, not an isolation guarantee — privileged containers or a compromised host can evade it. \textbf{Convention, not yet adopted:} CEP is designed and optional; both maturity gates (Section~\ref{subsec:params}) remain unmet, and the open problems of Section~\ref{subsec:open} — treasury governance, the zero-knowledge verification phase, and weak subjectivity for late-joining verifiers — are unsolved here. None of these is hidden; each bounds what the paper claims.

\subsection{Future work}

The near- and mid-term roadmap follows the open problems. On the governance track: build the Gate-2 prerequisites (a DID-bound relying-party registry with cluster attribution, an explicit \emph{convention-status} signal, multi-chain anchoring), then exercise the testnet ramp; specify the operator-signed override-record format that makes scope-(a) enforcement turn-key; and design the zero-knowledge verification phase and a weak-subjectivity reference point for late joiners. On the substrate track: adversarial-scale Sybil validation and trust-weight sensitivity analysis at larger N; expressive-policy support to close the AIP open item; and an independent third-party security audit. The governance items are the paper's centre of gravity; the substrate items keep the supporting evidence honest as it scales.

% ---------------------------------------------------------------------
\section{Conclusion}
\label{sec:conclusion}
% Material: v1.9 §8 reframed — optional convention + maturity signal +
% deployed advisory substrate + operator-local-available-today + RFC invitation.
% ---------------------------------------------------------------------

A network-wide trust signal that operators defer to must answer who decides it is ready — without that decision becoming a new central operator. This paper separated the question into two scopes. \textbf{Operator-local} enforcement — a relying party gating its own agent against its own mandate — needs no convention and is available on the deployed substrate today; it makes no claim over anyone else's agents. The \textbf{network-wide} scope is an \emph{optional} trust convention the Combined Enforcement Protocol governs: a five-condition maturity signal whose value is a deterministic function any party can recompute from anchored data, with keyed commitments reconciling permanence against the right to erasure, staged verification toward a zero-knowledge end-state, and cluster diversity proven from graph structure rather than declared categories. The design dissolves the oracle problem that conditions-based recognition usually reintroduces, and its concentration caps are stricter than the on-chain-governance norm. CEP compels no one: it offers a signal to coordinate on, never an obligation.

We have been deliberate about status. The layer runs \textbf{advisory} in the reference deployment; the convention is \textbf{designed} and optional, its maturity signal gated, not active; the mechanism is \textbf{demonstrable on testnet}; and the W3C VC + DID substrate — live since March 2026, anchored on Base L2 — is the supporting evidence that the substrate is real, with conformance reported honestly as four AIP features met and one open. Operator-independence here is a property of the maturity \emph{signal} — that no single party certifies the convention's readiness — not a claim to enforce on anyone, and we have drawn that line throughout so it cannot be mistaken for one.

The design is settled; whether the convention matures and is adopted is not, and that is by intent. We invite the critique a trust convention deserves before it is relied upon: independent recomputation of the maturity predicate, scrutiny of the parameters, and — most valuable — relying parties willing to join the testnet ramp and exercise the mechanism. The companion RFC opens that process. The infrastructure regulators and industry have converged on is implementable now; an optional, operator-independent way to recognise when a network-wide trust convention has matured is a design we are ready to have tested.
